\section{Professional benchmarks and economic interpretation}\label{sec:C24:professional}
Professional-task benchmarks attempt to move evaluation beyond short questions into deliverables with multiple constraints. Their relevance is substantial, but the mapping to economic impact requires care. A benchmark task can resemble professional work while omitting coordination, accountability, access friction, and consequences of failure. A successful artifact is evidence about the tested task, not a direct estimate of organizational value or labor substitution.

Agents' Last Exam introduces long-horizon professional tasks with verifiable outcomes across occupational categories~\cite{corpus260605405}. Its living-benchmark design requires exact task-pool and scoring snapshots for longitudinal comparison. An improved score can reflect a changed task mix or grader as well as a stronger system. Reports should retain the task version, evaluator identity, available tools, and artifact requirements rather than cite a moving leaderboard alone.

Economic interpretation also depends on selection. Benchmarks often choose tasks with accessible data and checkable outcomes. Work requiring tacit negotiation, physical access, or responsibility allocation may be underrepresented. This selection does not invalidate the benchmark; it limits the population to which results generalize. A transparent survey should state which professional activities the evaluation captures and which remain outside its observation boundary.

Task cost must include verification and remediation. Let gross production benefit be $V$, execution cost $C_e$, review cost $C_v$, and expected failure loss $L$. A deployment decision depends on $V-C_e-C_v-L$ under a specified workload. These quantities may be uncertain and correlated. Faster generation can increase review volume, while stronger checks can reduce failure loss at additional cost. No universal conversion from token throughput to economic value follows.

The proposed evaluation therefore combines benchmark tasks with deployment-specific shadow runs. Shadow execution observes what the agent would propose without granting consequential effects. It can reveal permission and evidence failures while preserving a human-controlled process. It does not fully measure autonomous deployment because humans and environments may behave differently when the agent's actions are real. The limitation should remain explicit when interpreting readiness.

